%%%%%%%%%%%%%%%%%%%%%%%%%%%%%%%%%%%%%%%%%%%%%%%%%%%%%%%%%%%%%%%%%%%%%%%%
%                                                                      %
%     File: Thesis_Resumo.tex                                          %
%     Tex Master: Thesis.tex                                           %
%                                                                      %
%     Author: Andre C. Marta                                           %
%     Last modified :  4 Mar 2024                                      %
%                                                                      %
%%%%%%%%%%%%%%%%%%%%%%%%%%%%%%%%%%%%%%%%%%%%%%%%%%%%%%%%%%%%%%%%%%%%%%%%

\section*{Resumo}

% Add entry in the table of contents as section
\addcontentsline{toc}{section}{Resumo}

Uma embarcação de superfície não tripulada (USV) e um veículo aéreo não tripulado (UAV) ligados por um cabo combinam a autonomia do USV com a visão aérea do UAV, que recebe energia e comunicação fiável através do cabo. O cabo, no entanto, acopla os dois veículos, que têm de se manter ao seu alcance respeitando cada um os seus limites, sob ondas e vento. Esta tese desenvolve o controlo cooperativo deste sistema, formado pelo USV WAM-V e pelo UAV x500. São derivados e validados num simulador modelos de ambos os veículos, bem como um modelo analítico da força de um cabo com folga que considera o seu peso. Cada veículo é comandado pelo seu próprio controlador preditivo não linear, e os dois controladores trocam as trajetórias que preveem, de modo que cada um antecipa o movimento do outro dentro do alcance do cabo. Um observador de perturbações compensa as forças não modeladas, como o vento. Numa simulação no mar, o controlador do UAV segue uma trajetória com menos de um quarto do erro do controlador de posição do PX4, e ambos os controladores mantêm os veículos dentro dos seus limites quando a referência não pode ser seguida. O UAV aterra no deck em movimento a poucos centímetros do centro da plataforma com ondas até $1{,}2\,\text{m}$ e ventos até $8\,\text{m/s}$, e em missões de inspeção de uma torre, perseguição de um alvo e recuperação do UAV, os dois veículos operam como um só sistema.

\vfill

\textbf{\Large Palavras-chave:} Cabo, Veículo Aéreo Não Tripulado, Embarcação de Superfície Não Tripulada, Controlo Cooperativo, Controlo Preditivo Não Linear

